%%% source
%%% book: zorich-1
%%% chapter: 1
%%% section: 2
%%% title: Множества и элементарные операции над множествами
%%% pages: 4-10
%%%

\section{Множества и элементарные операции над множествами}

\subsection{Понятие множества}

С конца XIX~--- начала XX столетия наиболее универсальным языком математики стал язык теории множеств. Это проявилось даже в одном из определений математики как науки, изучающей различные структуры (отношения) на множествах\footnotemark[1].

\footnotetext[1]{\emph{Бурбаки\,Н.} Архитектура математики. В кн.: \emph{Бурбаки\,Н.} Очерки по истории математики. М.: ИЛ, 1963.}

<<Под \emph{множеством} мы понимаем одно целое определённых, вполне различимых объектов нашей интуиции или нашей мысли>>~--- так описал понятие <<множество>> Георг Кантор\footnotemark[2], основатель теории множеств.

\footnotetext[2]{Г.\,Кантор (1845--1918)~--- немецкий математик, создатель теории бесконечных множеств и родоначальник теоретико-множественного языка в математике.}

Описание Кантора, разумеется, нельзя назвать определением, поскольку оно апеллирует к понятиям, быть может, более сложным (во всяком случае, не определённым ранее), чем само понятие множества. Цель этого описания~--- разъяснить понятие, связав его с другими.

Основные предпосылки канторовской (или, как условно говорят, <<наивной>>) теории множеств сводятся к следующему:

$1^\circ$ множество может состоять из любых различимых объектов;

$2^\circ$ множество однозначно определяется набором составляющих его объектов;

$3^\circ$ любое свойство определяет множество объектов, которые этим свойством обладают.

Если $x$~--- объект, $P$~--- свойство, $P(x)$~--- обозначение того, что $x$ обладает свойством $P$, то через $\{x \mid P(x)\}$ обозначают весь класс объектов, обладающих свойством $P$. Объекты, составляющие класс или множество, называют \emph{элементами} класса или множества.

Множество, состоящее из элементов $x_1, \ldots, x_n$, обычно обозначают как $\{x_1, \ldots, x_n\}$. Там, где это не вызывает недоразумения, для сокращения записи мы позволяем себе обозначать одноэлементное множество $\{a\}$ просто через~$a$.

Слова <<класс>>, <<семейство>>, <<совокупность>>, <<набор>> в наивной теории множеств употребляют как синонимы термина <<множество>>.

Следующие примеры демонстрируют применение этой терминологии:

множество букв <<а>> в слове <<я>>;

множество жён Адама;

набор из десяти цифр;

семейство бобовых;

множество песчинок на Земле;

совокупность точек плоскости, равноудалённых от двух данных её точек;

семейство множеств;

множество всех множеств.

Различие в возможной степени определённости задания множества наводит на мысль, что множество~--- не такое уж простое и безобидное понятие. И в самом деле, например, понятие множества всех множеств просто противоречиво.

$\blacktriangleleft$ Действительно, пусть для множества $M$ запись $P(M)$ означает, что $M$ не содержит себя в качестве своего элемента.

Рассмотрим класс $K = \{M \mid P(M)\}$ множеств, обладающих свойством $P$.

Если $K$~--- множество, то либо верно, что $P(K)$, либо верно, что $\neg P(K)$. Однако эта альтернатива для $K$ невозможна. Действительно, $P(K)$ невозможно, ибо из определения $K$ тогда бы следовало, что $K$ содержит $K$, т.\,е. что верно $\neg P(K)$; с другой стороны, $\neg P(K)$ тоже невозможно, поскольку это означает, что $K$ содержит $K$, а это противоречит определению $K$ как класса тех множеств, которые сами себя не содержат.

Следовательно, $K$~--- не множество. $\blacktriangleright$

Это классический парадокс Рассела\footnotemark[1], один из тех парадоксов, к которым приводит наивное представление о множестве.

\footnotetext[1]{Б.\,Рассел (1872--1970)~--- английский логик, философ, социолог и общественный деятель.}

В современной математической логике понятие множества подвергается (как мы видим, не без оснований) тщательному анализу. Однако в такой анализ мы углубляться не станем. Отметим только, что в существующих аксиоматических теориях множество определяется как математический объект, обладающий определённым набором свойств.

Описание этих свойств составляет аксиоматику. Ядром аксиоматики теории множеств является постулирование правил, по которым из множеств можно образовывать новые множества. В целом любая из существующих аксиоматик такова, что она, с одной стороны, избавляет от известных противоречий наивной теории, а с другой~--- обеспечивает свободу оперирования с конкретными множествами, возникающими в различных отделах математики, и в первую очередь именно в математическом анализе, понимаемом в широком смысле слова.

Ограничившись пока этими замечаниями относительно понятия множества, перейдём к описанию некоторых наиболее часто используемых в анализе свойств множеств.

Желающие подробнее ознакомиться с понятием множества могут просмотреть пункт~2 из \S\;4 настоящей главы или обратиться к специальной литературе.

\subsection{Отношение включения}

Как уже отмечалось, объекты, составляющие множество, принято называть \emph{элементами} этого множества. Мы будем стремиться обозначать множества прописными буквами латинского алфавита, а элементы множества~--- соответствующими строчными буквами.

Высказывание <<$x$ есть элемент множества $X$>> коротко обозначают символом
\[
x \in X \quad (\text{или } X \ni x),
\]
а его отрицание~--- символом
\[
x \notin X \quad (\text{или } X \not\ni x).
\]

В записи высказываний о множествах часто используются логические операторы $\exists$ (<<существует>> или <<найдётся>>) и $\forall$ (<<любой>> или <<для любого>>), называемые кванторами \emph{существования} и \emph{всеобщности} соответственно.

Например, запись $\forall\,x\;((x \in A) \Leftrightarrow (x \in B))$ означает, что для любого объекта $x$ соотношения $x \in A$ и $x \in B$ равносильны. Поскольку множество вполне определяется своими элементами, указанное высказывание принято обозначать короткой записью
\[
A = B,
\]
читаемой <<$A$ равно $B$>>, обозначающей совпадение множеств $A$ и $B$.

Таким образом, два множества \emph{равны}, когда они состоят из одних и тех же элементов.

Отрицание равенства обычно записывают в виде $A \neq B$.

Если любой элемент множества $A$ является элементом множества $B$, то пишут $A \subset B$ или $B \supset A$ и говорят, что множество $A$ является \emph{подмножеством} множества $B$, или что $B$ содержит $A$, или что $B$ включает в себя $A$. В связи с этим отношение $A \subset B$ между множествами $A$, $B$ называется \emph{отношением включения} (рис.\;1).

Итак,
\[
(A \subset B) := \forall\,x\;((x \in A) \Rightarrow (x \in B)).
\]

Если $A \subset B$ и $A \neq B$, то будем говорить, что включение $A \subset B$ \emph{строгое} или что $A$~--- \emph{собственное} подмножество $B$.

Используя приведённые определения, теперь можно заключить, что
\[
(A = B) \Leftrightarrow (A \subset B) \land (B \subset A).
\]

Если $M$~--- множество, то любое свойство $P$ выделяет в $M$ подмножество
\[
\{x \in M \mid P(x)\}
\]
тех элементов $M$, которые обладают этим свойством.

Например, очевидно, что
\[
M = \{x \in M \mid x \in M\}.
\]

С другой стороны, если в качестве $P$ взять свойство, которым не обладает ни один элемент множества $M$, например $P(x) := (x \neq x)$, то мы получим множество
\[
\varnothing = \{x \in M \mid x \neq x\},
\]
называемое \emph{пустым} подмножеством множества $M$.

\subsection{Простейшие операции над множествами}

Пусть $A$ и $B$~--- подмножества множества $M$.

\textbf{a.} \emph{Объединением} множеств $A$ и $B$ называется множество
\[
A \cup B := \{x \in M \mid (x \in A) \lor (x \in B)\},
\]
состоящее из тех и только тех элементов множества $M$, которые содержатся хотя бы в одном из множеств $A$, $B$ (рис.\;2).

\textbf{b.} \emph{Пересечением} множеств $A$ и $B$ называется множество
\[
A \cap B := \{x \in M \mid (x \in A) \land (x \in B)\},
\]
образованное теми и только теми элементами множества $M$, которые принадлежат одновременно $A$ и $B$ (рис.\;3).

\textbf{c.} \emph{Разностью} между множеством $A$ и множеством $B$ называется множество
\[
A \setminus B := \{x \in M \mid (x \in A) \land (x \notin B)\},
\]
состоящее из тех элементов множества $A$, которые не содержатся в множестве $B$ (рис.\;4).

Разность между множеством $M$ и содержащимся в нём подмножеством $A$ обычно называют \emph{дополнением} $A$ в $M$ и обозначают через $C_M A$ или $CA$, когда из контекста ясно, в каком множестве ищется дополнение к $A$ (рис.\;5).

\textbf{Пример.} В качестве иллюстрации взаимодействия введённых понятий проверим следующие соотношения (так называемые правила де Моргана\footnotemark[1]):

\footnotetext[1]{А.\,де Морган (1806--1871)~--- шотландский математик.}

\begin{align}
C_M(A \cup B) &= C_M A \cap C_M B, \label{eq:demorgan1} \\
C_M(A \cap B) &= C_M A \cup C_M B. \label{eq:demorgan2}
\end{align}

$\blacktriangleleft$ Докажем, например, первое из этих равенств:
\begin{align*}
(x \in C_M(A \cup B)) &\Rightarrow (x \notin (A \cup B)) \Rightarrow ((x \notin A) \land (x \notin B)) \Rightarrow \\
&\Rightarrow (x \in C_M A) \land (x \in C_M B) \Rightarrow (x \in (C_M A \cap C_M B)).
\end{align*}

Таким образом, установлено, что
\begin{equation}
C_M(A \cup B) \subset (C_M A \cap C_M B). \label{eq:dm-forward}
\end{equation}

С другой стороны,
\begin{align*}
(x \in (C_M A \cap C_M B)) &\Rightarrow ((x \in C_M A) \land (x \in C_M B)) \Rightarrow \\
&\Rightarrow ((x \notin A) \land (x \notin B)) \Rightarrow (x \notin (A \cup B)) \Rightarrow (x \in C_M(A \cup B)),
\end{align*}
т.\,е.
\begin{equation}
(C_M A \cap C_M B) \subset C_M(A \cup B). \label{eq:dm-backward}
\end{equation}

Из \eqref{eq:dm-forward} и \eqref{eq:dm-backward} следует \eqref{eq:demorgan1}. $\blacktriangleright$

\textbf{d.} \emph{Прямое (декартово) произведение множеств.} Для любых двух множеств $A$, $B$ можно образовать новое множество~--- пару $\{A, B\} = \{B, A\}$, элементами которого являются множества $A$ и $B$ и только они. Это множество состоит из двух элементов, если $A \neq B$, и из одного элемента, если $A = B$.

Указанное множество называют \emph{неупорядоченной парой} множеств $A$, $B$, в отличие от \emph{упорядоченной пары\/} $(A, B)$, в которой элементы $A$, $B$ наделены дополнительными признаками, выделяющими первый и второй элементы пары $\{A, B\}$. Равенство
\[
(A, B) = (C, D)
\]
упорядоченных пар по определению означает, что $A = C$ и $B = D$. В частности, если $A \neq B$, то $(A, B) \neq (B, A)$.

Пусть теперь $X$ и $Y$~--- произвольные множества. Множество
\[
X \times Y := \{(x, y) \mid (x \in X) \land (y \in Y)\},
\]
образованное всеми упорядоченными парами $(x, y)$, первый член которых есть элемент из $X$, а второй член~--- элемент из $Y$, называется \emph{прямым} или \emph{декартовым произведением множеств} $X$ и $Y$ (в таком порядке!).

Из определения прямого произведения и сделанных выше замечаний об упорядоченной паре явствует, что, вообще говоря, $X \times Y \neq Y \times X$. Равенство имеет место, лишь если $X = Y$. В последнем случае вместо $X \times X$ пишут коротко $X^2$.

Прямое произведение называют также декартовым произведением в честь Декарта\footnotemark[1], который независимо от Ферма\footnotemark[2] пришёл через систему координат к аналитическому языку геометрии. Известная всем система декартовых координат в плоскости превращает эту плоскость именно в прямое произведение двух числовых осей. На этом знакомом объекте наглядно проявляется зависимость декартова произведения от порядка сомножителей. Например, упорядоченным парам $(0, 1)$ и $(1, 0)$ отвечают различные точки плоскости.

\footnotetext[1]{Р.\,Декарт (1596--1650)~--- выдающийся французский философ, математик и физик, внёсший фундаментальный вклад в теорию научного мышления и познания.}

\footnotetext[2]{П.\,Ферма (1601--1665)~--- замечательный французский математик, юрист по специальности. Ферма стоял у истоков ряда областей современной математики: анализ, аналитическая геометрия, теория вероятностей, теория чисел.}

В упорядоченной паре $z = (x_1, x_2)$, являющейся элементом прямого произведения $Z = X_1 \times X_2$ множеств $X_1$ и $X_2$, элемент $x_1$ называется \emph{первой проекцией пары} $z$ и обозначается через $\operatorname{pr}_1 z$, а элемент $x_2$~--- \emph{второй проекцией пары} $z$ и обозначается через $\operatorname{pr}_2 z$.

Проекции упорядоченной пары по аналогии с терминологией аналитической геометрии часто называют (первой и второй) \emph{координатами пары}.

\subsection*{Упражнения}

В задачах 1, 2, 3 через $A$, $B$, $C$ обозначены подмножества некоторого множества $M$.

\textbf{1.} Проверьте соотношения

a) $(A \subset C) \land (B \subset C) \Leftrightarrow ((A \cup B) \subset C)$; \qquad b) $(C \subset A) \land (C \subset B) \Leftrightarrow (C \subset (A \cap B))$;

c) $C_M(C_M A) = A$; \qquad d) $(A \subset C_M B) \Leftrightarrow (B \subset C_M A)$;

e) $(A \subset B) \Leftrightarrow (C_M A \supset C_M B)$.

\textbf{2.} Покажите, что

a) $A \cup (B \cup C) = (A \cup B) \cup C =: A \cup B \cup C$; \quad b) $A \cap (B \cap C) = (A \cap B) \cap C =: A \cap B \cap C$;

c) $A \cap (B \cup C) = (A \cap B) \cup (A \cap C)$; \quad d) $A \cup (B \cap C) = (A \cup B) \cap (A \cup C)$.

\textbf{3.} Проверьте взаимосвязь (двойственность) операций объединения и пересечения:

a) $C_M(A \cup B) = C_M A \cap C_M B$; \qquad b) $C_M(A \cap B) = C_M A \cup C_M B$.

\textbf{4.} Проиллюстрируйте геометрически декартово произведение

a) двух отрезков (прямоугольник); \qquad b) двух прямых (плоскость);

c) прямой и окружности (цилиндрическая поверхность);

d) прямой и круга (цилиндр); \qquad e) двух окружностей (тор);

f) окружности и круга (полноторие).

\textbf{5.} Множество $\Delta = \{(x_1, x_2) \in X^2 \mid x_1 = x_2\}$ называется \emph{диагональю декартова квадрата} $X^2$ множества $X$.

Проиллюстрируйте геометрически диагонали множеств, полученных в пунктах a), b), e) задачи~4.

\textbf{6.} Покажите, что

a) $(X \times Y = \varnothing) \Leftrightarrow (X = \varnothing) \lor (Y = \varnothing)$,

а если $X \times Y \neq \varnothing$, то

b) $(A \times B \subset X \times Y) \Leftrightarrow (A \subset X) \land (B \subset Y)$, \quad c) $(X \times Y) \cup (Z \times Y) = (X \cup Z) \times Y$,

d) $(X \times Y) \cap (X' \times Y') = (X \cap X') \times (Y \cap Y')$.

Здесь $\varnothing$~--- символ пустого множества, т.\,е. множества, не содержащего элементов.

\textbf{7.} Сравнив соотношения задачи~3 с соотношениями a), b) из упражнения~2 к \S\;1, установите соответствие между логическими операциями $\neg$, $\land$, $\lor$ на высказываниях и операциями $C$, $\cap$, $\cup$ на множествах.
